\documentclass[aspectratio=169]{beamer}

\usetheme{Madrid}
\usecolortheme{seahorse}
\setbeamertemplate{navigation symbols}{}
\setbeamertemplate{footline}[frame number]

% To print slides WITH speaker notes, uncomment:
% \setbeameroption{show notes}

\usepackage{amsmath, amssymb}
\usepackage{graphicx}

% Placeholder helper: shows image if present, else a labeled box.
\newcommand{\circuitfig}[3][\linewidth]{%
  \IfFileExists{#2}{\includegraphics[width=#1]{#2}}{%
    \fbox{\parbox[c][3.2cm][c]{0.9\linewidth}{\centering \texttt{#2}\\[2pt]\small #3}}}%
}

\title[QFT, QPE, and Shor's Algorithm]{The Quantum Fourier Transform,\\ Phase Estimation, and Shor's Algorithm}
\author{M.~Thompson}
\institute{Math 481 --- Complex Analysis}
\date{\today}

\begin{document}

\begin{frame}
  \titlepage
\end{frame}
\note{
The punchline of this talk: quantum computers can factor big numbers fast, and
the engine that makes it work is the Fourier transform --- the same one from
this course, applied to a quantum state. Four steps: DFT, QFT, phase
estimation, Shor's. Don't memorize matrices; track what each piece is doing.
}

%----------------------------------------------------------------------
\begin{frame}{Roadmap}
  \vfill
  \begin{center}
  {\Huge DFT $\rightarrow$ QFT $\rightarrow$ QPE $\rightarrow$ Shor's}
  \end{center}
  \vfill
  \begin{center}
  Goal: factor a large integer $N$ in time polynomial in $\log N$.
  \end{center}
  \vfill
\end{frame}
\note{
Each arrow means ``the previous thing, used as a tool.'' The QFT is the DFT in
disguise. QPE is a clever use of the QFT. Shor's is QPE pointed at one special
matrix. If the audience remembers one thing, it's this chain --- the rest of the
talk just fills in the arrows. Classically the best factoring algorithms are
sub-exponential; Shor's is polynomial in the number of digits, which is why it
breaks RSA in principle.
}

%----------------------------------------------------------------------
\begin{frame}{The DFT: a change of basis from roots of unity}
  \[
    \hat{x}_k = \frac{1}{\sqrt{N}}\sum_{l=0}^{N-1} x_l\, \omega^{kl},
    \qquad \omega = e^{-2\pi i/N}
  \]
  \begin{block}{Orthogonality of roots of unity}
  \[
    \frac{1}{N}\sum_{k=0}^{N-1}\omega^{mk} =
    \begin{cases} 1 & N \mid m \\ 0 & \text{else} \end{cases}
  \]
  Proof for $N \nmid m$: with $r = \omega^m \neq 1$, geometric series gives
  \(
    \frac{1}{N}\cdot\frac{r^N - 1}{r-1} = \frac{1}{N}\cdot\frac{e^{-2\pi i m}-1}{r-1} = 0.
  \)
  \end{block}
\end{frame}
\note{
Remind them: the DFT takes a list of numbers to a new list; each output is a
weighted sum where the weights are roots of unity --- points evenly spaced on
the unit circle. The block is the one proof worth showing in full, and it's
pure complex analysis: if N doesn't divide m, the powers of omega-to-the-m walk
around the circle and the geometric series sums to zero, because $r^N = 1$
exactly. All the roots cancel. This cancellation is the reason the Fourier
basis is orthonormal, hence the reason the DFT matrix is unitary, hence ---
much later --- the reason the quantum circuit is even allowed to exist. One
slide of complex analysis powers the whole talk.
}

%----------------------------------------------------------------------
\begin{frame}{Qubits}
  \begin{columns}
  \begin{column}{0.5\textwidth}
    \[
      \vert\psi\rangle = \alpha\vert 0\rangle + \beta\vert 1\rangle,
      \qquad |\alpha|^2 + |\beta|^2 = 1
    \]
    \[
      n \text{ qubits } \longleftrightarrow\ \mathbb{C}^{2^n} = (\mathbb{C}^2)^{\otimes n}
    \]
  \end{column}
  \begin{column}{0.45\textwidth}
    \begin{center}
      \circuitfig[0.62\linewidth]{Bloch_sphere.png}{Bloch sphere}
    \end{center}
  \end{column}
  \end{columns}
\end{frame}
\note{
Skip slide --- say ``we did this in class.'' A qubit is a unit vector in C2;
n qubits live in a $2^n$-dimensional space via tensor products. The picture is
the Bloch sphere: after modding out the global phase, a single qubit state is a
point on an ordinary sphere --- and for us complex analysts, that sphere is the
Riemann sphere: the ratio beta over alpha lives in C union infinity, and
stereographic projection puts it on S2. North pole is 0, south pole is 1. Gates
act as rotations of this sphere. Move on quickly.
}

%----------------------------------------------------------------------
\begin{frame}{The gate set}

\begin{columns}
\begin{column}{\textwidth}
\[
H = \frac{1}{\sqrt{2}}
\begin{bmatrix}
1 & 1\\
1 & -1
\end{bmatrix}
\qquad
R_k =
\begin{bmatrix}
1 & 0\\
0 & e^{2\pi i/2^k}
\end{bmatrix}
\]
\end{column}
\end{columns}

\vspace{0.5em}

\begin{columns}
\begin{column}{\textwidth}
\[
\mathrm{SWAP} =
\begin{bmatrix}
1&0&0&0\\
0&0&1&0\\
0&1&0&0\\
0&0&0&1
\end{bmatrix}
\quad
C(R_k) =
\begin{bmatrix}
1&0&0&0\\
0&1&0&0\\
0&0&1&0\\
0&0&0&e^{2\pi i/2^k}
\end{bmatrix}
\]
\end{column}
\end{columns}

\end{frame}
\note{
Four gates, the whole toolkit. Hadamard turns a definite 0 or 1 into an even
mix --- it creates superpositions. R-k leaves 0 alone and spins the 1-component
by a fixed angle: it writes a phase. SWAP exchanges two wires. Controlled R-k
fires the phase only when the control qubit is 1 --- point at the bottom-right
entry, that single $e^{2\pi i/2^k}$ is its only action. On the Bloch sphere
picture from the previous slide, the R-k gates are rotations about the vertical
axis. Everything in the QFT circuit is built from just these.
}

%----------------------------------------------------------------------
\begin{frame}{QFT: the DFT factorizes over qubits}
  \[
    F_{2^n}\vert j_1 \cdots j_n\rangle =
    \frac{1}{\sqrt{2^n}}
    \left(\vert 0\rangle + e^{2\pi i\, 0.j_n}\vert 1\rangle\right)
    \otimes \cdots \otimes
    \left(\vert 0\rangle + e^{2\pi i\, 0.j_1 \cdots j_n}\vert 1\rangle\right)
  \]
  \vfill
  \begin{center}
    One qubit per factor \ $\Rightarrow$ \ $O(n^2)$ gates instead of $O(2^n)$.
  \end{center}
\end{frame}
\note{
The QFT is the DFT applied to amplitudes of a quantum state. Naively that's a
$2^n$ by $2^n$ matrix --- hopeless. But when N is a power of two, the output
miraculously splits into a tensor product of single-qubit states; that's the
formula, with $0.j \ldots$ meaning a binary fraction. Each factor is one qubit
holding one phase, built by a Hadamard plus a few controlled rotations. So an
exponential-size transform collapses to about n-squared gates. Plain-language
version to say out loud: the QFT takes information stored in amplitudes and
re-writes it as phases --- angles --- spread across the qubits. Hold onto that;
phases are exactly what the next tool reads out.
}

%----------------------------------------------------------------------
\begin{frame}{QFT circuit}
  \begin{center}
    \circuitfig{QFT_circuit.png}{QFT circuit}
  \end{center}
\end{frame}
\note{
Walk it left to right: each wire gets a Hadamard, then controlled rotations ---
the S, T, R4, R5 boxes are R-k gates --- reach over from the lower qubits to
stamp in smaller and smaller phase corrections, the binary-fraction digits from
the last formula. The crossings at the far right are SWAPs putting the qubits
back in order. Don't trace every gate; the point is the circuit is short and
regular --- that regularity is the n-squared efficiency made visible.
}

%----------------------------------------------------------------------
\begin{frame}{QPE: eigenvalues of unitaries are phases}
  \begin{columns}
  \begin{column}{0.55\textwidth}
    If $U\vert\psi\rangle = \lambda\vert\psi\rangle$ with $U$ unitary, then
    \[
      |\lambda| = 1, \qquad \lambda = e^{2\pi i\theta},\ \ \theta \in [0,1).
    \]
    \begin{block}{Proof}
    \(
      \|\vert\psi\rangle\| = \|U\vert\psi\rangle\| = |\lambda| \|\vert\psi\rangle\|
      \Rightarrow |\lambda| = 1.
    \)
    \end{block}
    Goal of QPE: find $\theta$.
  \end{column}
  \begin{column}{0.4\textwidth}
    \begin{center}
      \circuitfig[0.75\linewidth]{Phase_UnitCircle.PNG}{eigenvalue on the unit circle}
    \end{center}
  \end{column}
  \end{columns}
\end{frame}
\note{
Second tool. A unitary preserves norms, so the one-line proof on the slide
shows every eigenvalue has modulus one --- eigenvalues live ON the unit circle,
as in the picture. So an eigenvalue is nothing but an angle theta. Phase
estimation asks: given the matrix and an eigenvector, find that angle. The
catch: applying U to its eigenvector only multiplies by a global phase, and
global phases are invisible to measurement --- probabilities only see absolute
values. So we need a trick to move theta somewhere measurable. That trick is
phase kickback, next slide.
}

%----------------------------------------------------------------------
\begin{frame}{Phase kickback}
  \[
    \tfrac{1}{\sqrt2}(\vert 0\rangle + \vert 1\rangle)\otimes\vert\psi\rangle
    \ \xrightarrow{\ C(U)\ }\
    \tfrac{1}{\sqrt2}\bigl(\vert 0\rangle + e^{2\pi i\theta}\vert 1\rangle\bigr)\otimes\vert\psi\rangle
  \]
  \vfill
  \begin{center}
    When we do a controlled gate, controlled by a qubit that has been put into superposition
    by a Hadamard gate, the control basically gets "reversed."
    
    \vspace{15pt}
    
    The target is unchanged, the phase lands on the \emph{control}.
  \end{center}
\end{frame}
\note{
The cleverest move in the talk; slow down here. Normally the control decides
and the target gets acted on. We flip it: put the CONTROL in superposition and
feed the TARGET an eigenvector. On the 0-branch nothing happens; on the
1-branch U hits psi --- but U on its eigenvector is just multiplication by
$e^{2\pi i \theta}$, a number, which slides off psi and attaches to the
1-branch of the control. End result: psi is untouched, and the phase has jumped
``backwards'' onto the control qubit. The gate ran in reverse, in a sense ---
the target reprogrammed the control. That rescues the invisible phase: it now
sits on a qubit we can play with.
}

%----------------------------------------------------------------------
\begin{frame}{One ancilla: estimate $\theta$ by sampling}
  \begin{center}
    \circuitfig[0.55\linewidth]{1qubitancilla_QPE.PNG}{1-ancilla QPE circuit}
  \end{center}
  \[
    P(0) = \cos^2(\pi\theta), \qquad P(1) = \sin^2(\pi\theta)
  \]
\end{frame}
\note{
This circuit is the smallest possible phase estimator: Hadamard, controlled-U,
Hadamard, measure. The second Hadamard makes the 0 and 1 branches interfere,
and the interference angle is theta --- giving the biased coin on the slide:
probability cos-squared of pi theta for 0, sin-squared for 1. So run it a
thousand times, count zeros, solve backwards for theta --- a Monte Carlo
estimate, like estimating a coin's bias by flipping it repeatedly. It works but
it's slow: lots of shots for a few digits. Next slide: get all the digits in
ONE shot.
}

%----------------------------------------------------------------------
\begin{frame}{One ancilla: where $\cos^2(\pi\theta)$ comes from}
  After kickback, apply the second Hadamard to $\tfrac{1}{\sqrt2}\bigl(\vert 0\rangle + e^{2\pi i\theta}\vert 1\rangle\bigr)$:
  \[
    H\vert 0\rangle = \tfrac{1}{\sqrt2}(\vert 0\rangle + \vert 1\rangle),
    \qquad
    H\vert 1\rangle = \tfrac{1}{\sqrt2}(\vert 0\rangle - \vert 1\rangle)
  \]
  \[
    \tfrac{1}{\sqrt2}\bigl(\vert 0\rangle + e^{2\pi i\theta}\vert 1\rangle\bigr)
    \ \xrightarrow{\ H\ }\
    \tfrac{1}{2}\bigl(1 + e^{2\pi i\theta}\bigr)\vert 0\rangle
    + \tfrac{1}{2}\bigl(1 - e^{2\pi i\theta}\bigr)\vert 1\rangle
  \]
  \vfill
  Born rule, using $|1 \pm e^{2\pi i\theta}|^2 = 2 \pm 2\cos(2\pi\theta)$:
  \[
    P(0) = \tfrac{1}{4}\bigl|1 + e^{2\pi i\theta}\bigr|^2 = \cos^2(\pi\theta),
    \qquad
    P(1) = \tfrac{1}{4}\bigl|1 - e^{2\pi i\theta}\bigr|^2 = \sin^2(\pi\theta)
  \]
\end{frame}

%----------------------------------------------------------------------
\begin{frame}{Many ancillas: the kickbacks applied to the ancilla are effectively a QFT}
  \[
    \tfrac{1}{\sqrt{2^t}}
    \left(\vert 0\rangle + e^{2\pi i\, 2^{t-1}\theta}\vert 1\rangle\right)
    \otimes \cdots \otimes
    \left(\vert 0\rangle + e^{2\pi i\, \theta}\vert 1\rangle\right)
    \;=\; F_{2^t}\vert\tilde\theta\rangle
  \]
  \vfill
  \begin{center}
    Apply $F_{2^t}^{-1}$, measure: $\theta$ in binary, one shot.
  \end{center}
\end{frame}
\note{
Use t ancillas; the k-th one controls U to the power 2-to-the-k, so it kicks
back theta doubled k times --- each ancilla captures a different binary digit
of theta. Now compare the resulting product state with the QFT factorization
from earlier: identical. The pile of kicked-back phases IS the Fourier
transform of theta written in binary. Theta is hiding in Fourier space --- and
we own the machine that inverts that. Run the inverse QFT, measure, and the
binary digits of theta drop out directly. One shot, no statistics; more
ancillas means more digits. THIS is why we built the QFT first: it's the
read-out device for phases.
}

%----------------------------------------------------------------------
\begin{frame}{QPE circuit}
  \begin{center}
    \circuitfig{QPE_circuit.png}{QPE circuit}
  \end{center}
\end{frame}
\note{
The whole algorithm in one picture, three blocks: Hadamards make the
superposition; the controlled powers of U do the kickback, loading theta into
Fourier space; the inverse-QFT box decodes it and we measure the binary answer.
Bottom wire is the eigenvector, which just sits there donating its phase ---
note it comes out unchanged.
}

%----------------------------------------------------------------------
\begin{frame}{Shor's: the number theory}
  Order of $a$: smallest $r$ with $a^r \equiv 1 \pmod N$. \quad Let $b = a^{r/2}$ ($r$ even):
  \[
    (b-1)(b+1) = b^2 - 1 \equiv 0 \pmod{N}
  \]
  \vfill
  If $b \not\equiv \pm 1$: $N$ divides the product but neither factor
  $\Rightarrow$ its primes split between them
  \[
    \Rightarrow\ \gcd(b-1,\ N)\ \text{is a nontrivial factor of}\ N.
  \]
\end{frame}
\note{
The classical half. Pick a coprime to N; its powers cycle with period r, the
order. If r is even, set b = a to the r-over-2, so b squared is 1 mod N, and
b-squared-minus-1 factors as (b-1)(b+1). N divides that product. Now the key
step --- say it slowly: provided b isn't congruent to plus or minus 1, N divides
NEITHER factor alone. The only way to divide a product but neither piece is for
N's primes to SPLIT --- some land in (b-1), the rest in (b+1). So gcd(b-1, N)
scoops up exactly the primes that landed on the left: it can't be 1 (some are
there) and can't be N (not all are), so it's an honest factor. And gcds are
cheap --- Euclid's algorithm. Random a works with probability at least one half.
Conclusion: factoring reduces to finding r.
}

%----------------------------------------------------------------------
\begin{frame}{Shor's: order-finding is QPE}
  \[
    U\vert x\rangle = \vert ax \bmod N\rangle
    \qquad\Longrightarrow\qquad
    \text{eigenvalues } e^{2\pi i\, s/r}
  \]
  \vfill
  \begin{center}
    QPE measures $\phi \approx s/r$
    \ $\rightarrow$ continued fractions give $r$
    \ $\rightarrow$ $\gcd(a^{r/2}\pm 1,\, N)$.
  \end{center}
  \vfill
  \begin{center}
  \small Example: $N=15$, $a=7$: $r=4$, $7^2 \equiv 4$, $\gcd(3,15)=3$, $\gcd(5,15)=5$.
  \end{center}
\end{frame}
\note{
The quantum half: choose U to be multiplication by a, mod N. It's unitary
because multiplication by an invertible element just permutes the basis states.
And because applying it r times returns you to start, its eigenvalues are
exactly the r-th roots of unity, e to the 2 pi i s over r --- the same circle
points from slide one. So QPE reads off a phase s over r, and the denominator
hands us the order. Pipeline: measure a decimal close to s/r, clean it into an
exact fraction with continued fractions, read r, finish with two gcds. Worked
example: N = 15, a = 7. The order is 4 since 7-to-the-4th is 2401, which is 1
mod 15. Then b = 49 = 4 mod 15, and gcd(3,15) = 3, gcd(5,15) = 5: out pops
15 = 3 times 5. One detail I'm sweeping under the rug: we can't prepare an
eigenvector since they depend on r --- but the state ``1'' happens to be an
equal superposition of all of them, and QPE handles the superposition by
linearity, returning a random s over r. That's in the paper if anyone asks.
}

%----------------------------------------------------------------------
\begin{frame}{Takeaways}
  \begin{center}
  \Large
  roots of unity $\rightarrow$ QFT $\rightarrow$ QPE $\rightarrow$ Shor's
  \end{center}
  \vfill
  \begin{itemize}
    \item QFT: amplitudes $\to$ phases, in $O(n^2)$ gates.
    \item QPE: kick an eigenphase onto ancillas, decode with $F^{-1}$.
    \item Shor's: factoring $=$ order-finding $=$ QPE on $\vert x\rangle \mapsto \vert ax \bmod N\rangle$.
  \end{itemize}
\end{frame}
\note{
Close the loop by running the chain backwards: Shor's needs the order;
order-finding is phase estimation; phase estimation reads phases with the
inverse QFT; and the QFT exists because roots of unity are orthogonal --- the
complex analysis we started with. Honest caveat to end on: today's hardware
doesn't have nearly enough error-corrected logical qubits to factor anything
RSA-sized, so this is a theorem about the future. Stop and take questions.
}

\end{document}